
\section{Seleccionar significa traducir una función a restricciones}
Antes de buscar una aleación se escribe la función: transmitir fuerza, contener presión, mantener una holgura, resistir una temperatura o proteger un sustrato. Se identifican restricciones que no pueden negociarse, variables que el diseñador puede cambiar y un objetivo medible. La menor densidad solo es un objetivo indirecto: lo que interesa puede ser la masa de una pieza que satisfaga todas las restricciones.

Para una barra a tracción de longitud fija, carga $F$ y tensión admisible académica $\sigma_{adm}$, el área mínima resistente es $F/\sigma_{adm}$ y la masa mínima es $\rho FL/\sigma_{adm}$. Bajo estas hipótesis interesa una razón $\sigma_{adm}/\rho$ alta. Si gobierna el alargamiento máximo, $\Delta L=FL/(EA)$ conduce a una masa proporcional a $\rho/E$. Una clasificación por resistencia específica puede cambiar al imponer rigidez. Si gobierna pandeo, la forma de sección y la longitud entran de otro modo; no se reutiliza automáticamente el mismo índice.

La tensión admisible del ejercicio no debe confundirse con un valor de catálogo sin factores, estadística ni condiciones. En un motor, además, las propiedades pueden depender de temperatura, tiempo, dirección y estado superficial. Un material más resistente puede exigir un tratamiento más complejo o ser menos tolerante a un defecto. La selección final necesita comprobar el mecanismo que realmente limita la pieza.
\section{Dispersión, escala y trazabilidad del dato}
Una probeta pulida y pequeña no reproduce todas las condiciones de una pieza con taladros, uniones y una superficie expuesta a ambiente. El volumen sometido a tensiones elevadas, el tamaño de defecto, la orientación del producto y la secuencia de fabricación pueden modificar la población relevante. No se obtiene una propiedad de diseño añadiendo simplemente la palabra «aeronáutico» a un valor medio.

Una ficha útil identifica composición o especificación, forma de producto, dirección, tratamiento, temperatura, método de ensayo y estadística del dato. También distingue valor nominal, mínimo especificado y admisible de diseño. Si faltan esas condiciones, la comparación puede servir para preselección, pero no para justificar resistencia o vida. El historial de una pieza añade otra capa: un mismo número de aleación no demuestra que dos piezas hayan vivido la misma exposición térmica o corrosiva.
\repaso{Dos materiales tienen igual módulo específico. ¿Pueden intercambiarse en un disco?}{No. Esa igualdad solo compara una combinación de dos propiedades. El disco también exige resistencia, fatiga, tenacidad, comportamiento térmico, procesabilidad, control de defectos y datos de diseño aplicables. La geometría y las condiciones de servicio deben permanecer en el análisis.}
